\documentclass[ ../main.tex]{subfiles}
\providecommand{\mainx}{..}
\begin{document}
\section{Conclusion}
\label{sec:conclusion}

We have defined the \emph{Bernoulli map}, a probabilistic function $\APFun{f}[\fprate][\fnrate] \colon \Set{X} \to \Maybe{\Set{Y}}$ that approximates an objective map $\Fun{f} \colon \Set{X} \to \Set{Y}$ with per-element errors modeled as independent Bernoulli trials. This generalizes the Bernoulli set model~\cite{aset} from indicator functions $\Set{X} \to \BitSet$ to arbitrary codomains.

The central theoretical result is the composition theorem (\cref{sec:composition}), which shows that the composition of independent Bernoulli maps is itself a Bernoulli map and derives the resulting error rates in closed form. The false negative rate of a composition satisfies $\fnrate_{g \circ f} = 1 - (1-\fnrate_1)(1-\fnrate_2)$, while the false positive rate depends on the interaction between the false positive value distribution of the first map and the domain of the second.

We also showed that the model specializes to several important cases: the Bernoulli set (indicator functions), the binary symmetric and asymmetric channels (Bernoulli Booleans), and approximate constants (maps from the unit type).

The entropy map (\cref{sec:entropy-maps}) provides a space-optimal realization of the first-order Bernoulli map model, and the singular hash map (\cref{sec:shm}) is a concrete construction that achieves the information-theoretic lower bound of $-\log_2 \fprate + \mu$ bits per element asymptotically. Trapdoor entropy maps extend this to privacy-preserving settings relevant to encrypted search~\cite{songWagnerPerrig2000,curtmola2006}.

\paragraph{Two perspectives, one theorem.}
The function-space viewpoint developed here takes a \emph{denotational} approach: every Bernoulli approximation is a latent function $\Fun{f'} \colon \Set{X} \to \Set{Y}$ approximated per-input, and the Kronecker factorization of the joint confusion matrix is the structural theorem.
A companion paper on random approximate algebraic types takes the \emph{constructive} counterpart: approximate types are built from Bool using sum, product, and exponential type constructors.
Product types preserve the Kronecker factorization (parameters grow additively); sum types break it (tag errors couple the components); and the exponential type $\Set{A} \to \Set{B}$ is precisely the function space studied here.
The two perspectives converge at the exponential type, arriving at the same factorization theorem from opposite directions.
Bernoulli sets are the special case $\Set{Y} = \BitSet$; the optional type $\Set{A} + \Unit$ (a sum) recovers Bernoulli set membership as a type-theoretic construction.

\paragraph{Future directions.}
Monte Carlo validation of the composition error bounds, particularly the dependence on the parameter $\rho$ governing false-positive overlap, remains to be developed.
The relational algebra operators (join, project, select) produce non-trivial error propagation when applied to approximate relations (sets of tuples), and merit further investigation.

\end{document}
